\section{Independent finite enclosures for the original policy catalogue}
\label{supp:independentvalidation119}

This section supplies the mathematical justification for validation D.  It is
a finite verification procedure for the original investment objective, using
classical convexity and the AVaR threshold representation
\citep{rockafellaruryasev2002,boyd2004}.  It is not a new optimization theorem,
an independent population experiment, or a claim of neural exclusivity.
The implementation reconstructs the original economic primitives separately
from the production terminal compiler.  Its mathematical guarantee still
uses the original domain, continuation-convexity, and implemented-terminal-policy
bounds; software independence does not remove these assumptions.

\subsection{Original objective and precise validation statement}
Fix $d=2$, $X=[0,1]^2$, $R=1/32$, $\beta=15/16$, and
$\alpha\in\{1/2,3/4\}$.  Put $t=1-\alpha$ and let $P$ be the
uniform law on $[-R,R]$.  Indices below are cyclic.  The original primitives are
\begin{align*}
 A(x)&=\frac18+\frac{x_1+x_2}{16},\qquad h(a)=a^2+4a^4,\\
 F_j(x,a,z)&=\frac1{16}+\frac{x_j}{2}+\frac{x_{j+1}}8\\
       &\quad+\frac{x_j(1-x_{j+1})}{16}+b_ja+s_jz,\\
 b&=(1/2,1/4),\quad s=(1,-1),\\
 S(x)&=\sum_{j=1}^2(x_j-5/8)^2
       +\frac18\sum_{j=1}^2(x_j-x_{j+1})^2\\
       &\quad+2(1/2-x_1-x_2)_+^2,\\
 S_T(x)&=2\sum_{j=1}^2(x_j-5/8)^2
       +\frac18\sum_{j=1}^2(x_j-x_{j+1})^2\\
       &\quad+2(1/2-x_1-x_2)_+^2.
\end{align*}
For bounded $w$, write
\[
 \rho_\alpha(w)=\inf_{\eta\in\mathbb R}
       \left\{\eta+\frac1t\mathbb E_P(w-\eta)_+\right\}.
\]
The last-date optimal value and the two root fixed-action objectives are
\begin{align*}
 G_y(v)&=S(y)+h(v)+\beta\rho_\alpha(S_T(F(y,v,\cdot))),\\
 V_1(y)&=\min_{0\le v\le A(y)}G_y(v),\\
 Q_x(a)&=S(x)+h(a)+\beta\rho_\alpha(V_1(F(x,a,\cdot))),\\
 J_x(a)&=S(x)+h(a)+\beta\rho_\alpha(J_1(F(x,a,\cdot))).
\end{align*}
Here $J_1$ is the unchanged implemented terminal policy, not a newly
optimized or independently reconstructed controller.  The original results
in Proposition~\ref{prop:investmentfoundations119} establish domain
invariance, existence, and strong state convexity of $V_1$ with modulus
$(4-5\beta/4)/2$.  The implemented terminal guarantee is
\begin{equation}\label{eq:independentterminalallowance119}
  V_1(y)\le J_1(y)\le V_1(y)+\varepsilon_1,
  \qquad y\in X,\qquad \varepsilon_1=2^{-20}.
\end{equation}
The validation below relies on this proved uniform guarantee and makes no
convexity assumption about $J_1$.

\begin{proposition}[Finite original-policy validation enclosure]
\label{prop:independentvalidation119}
Fix an original state $x\in X$ and feasible action $a\in[0,A(x)]$.
Let $M\ge1$ be an integer and put
\[
 z_k=-R+\frac{Rk}{M},\qquad y_k=F(x,a,z_k),
       \qquad k=0,\ldots,2M.
\]
At each $y_k$, choose any feasible $v_k$ and an exactly maximizing AVaR
density $q_k$ for $S_T(F(y_k,v_k,\cdot))$.  Define
\begin{align*}
 g_k&=h'(v_k)+\beta\mathbb E_P
                [q_k\,\partial_v S_T(F(y_k,v_k,Z))],\\
 \delta_k&=\operatorname{proj}_{[-v_k,A(y_k)-v_k]}(-g_k/2),\\
 \ell_k&=G_{y_k}(v_k)+g_k\delta_k+\delta_k^2,
       &u_k&=G_{y_k}(v_k).
\end{align*}
These are valid bounds $\ell_k\le V_1(y_k)\le u_k$, including
boundary actions.  For an $n$-vector $w$, let $\rho_\alpha^{[n]}(w)$
denote AVaR under equal atom weights $1/n$.  Form the vectors
\[
 \ell^{\mathrm{mid}}=(\ell_1,\ell_3,\ldots,\ell_{2M-1}),
 \qquad
 u^{\mathrm{end}}=(u_0,u_2,u_2,u_4,\ldots,u_{2M-2},u_{2M}).
\]
Thus each cell contributes both endpoints to $u^{\mathrm{end}}$;
its interior endpoints occur twice.  With $c=S(x)+h(a)$, define
\begin{align}
 L_Q&=c+\beta\rho_\alpha^{[M]}(\ell^{\mathrm{mid}}),&
 U_Q&=c+\beta\rho_\alpha^{[2M]}(u^{\mathrm{end}}).
 \label{eq:independentQ119}
\end{align}
Then the original objectives satisfy
\begin{equation}\label{eq:independentJ119}
 L_Q\le Q_x(a)\le U_Q,
 \qquad L_Q\le J_x(a)\le U_Q+\beta\varepsilon_1.
\end{equation}
Every quantity in the construction is rational when the original state,
action, and primitive-node actions are rational.  A finite terminal search
affects the widths, not the validity, of these inequalities.
\end{proposition}

\subsection{Proof: exact original primitives and boundary-safe node bounds}
\begin{proof}
We first verify that all evaluations concern the original feasible economy.
For $p,q\in[0,1]$, the state part
$p/2+q/8+p(1-q)/16$ is increasing in both arguments and lies in
$[0,5/8]$.  Since $A(x)\in[1/8,1/4]$, $0\le b_ja\le1/8$,
and $|s_jz|\le1/32$, every successor belongs to
$[1/32,27/32]^2\subset X$.  This applies both to a root action and to
each feasible terminal action.  The continuous terminal objective attains its
minimum on the compact feasible interval.

Fix a terminal state $y$ and feasible $v$.  Write
$F(y,v,z)=(w_1+z,w_2-z)$.  The aggregate-state squared hinge in
$S_T$ is constant in $z$ because the two shock loadings sum to zero.
The first two squared deviations contribute $4z^2$ and the cyclic
coupling term contributes $z^2$.  Consequently
\[
 p(z):=S_T(F(y,v,z))=p_0+p_1z+5z^2.
\]
This is a primitive identity.  Its evaluation at $-R,0,R$ determines
the polynomial exactly; three-point interpolation would not identify a
general nonlinear function.  At rational inputs, all three coefficients
are rational.

Here is an explicit exact maximizing density.  Set $L=2R\alpha$ and
\[
 u=\operatorname{proj}_{[-R,R-L]}
             \left(-\frac{p_1}{2p_2}-\frac L2\right),
 \quad p_2=5,\qquad
 q(z)=t^{-1}\boldsymbol1_{[-R,u]\,\cup\,[u+L,R]}(z).
\]
The central interval has length $L$, so $0\le q\le t^{-1}$ and
$\mathbb E_Pq=1$.  To verify maximality directly, let
$\eta=\max\{p(u),p(u+L)\}$.  An interior cut is centered on the
quadratic's vertex, so its endpoint values coincide and $p\le\eta$
inside the central interval and $p\ge\eta$ on its two tails.  At the
left boundary the clipped vertex lies weakly to the left of the central
midpoint, $p(u+L)\ge p(u)$, and the only nonempty tail is to the
right of $u+L$; it again has $p\ge\eta$.  The right-boundary case
is symmetric.  Thus, up to endpoints of probability zero,
\[
 \eta+t^{-1}\mathbb E_P(p-\eta)_+=\mathbb E_P(qp).
\]
For any feasible density $\widetilde q$ and any threshold $e$,
\[
 \mathbb E_P(\widetilde q p)
 =e+\mathbb E_P[\widetilde q(p-e)]
 \le e+t^{-1}\mathbb E_P(p-e)_+.
\]
Taking an infimum in $e$ proves weak duality; the displayed equality
proves that $q$ attains the original AVaR.  Integrating the quadratic
over its two rational tail intervals gives the exact rational risk.
This proof includes clipped cuts and does not differentiate a moving cut.

The function $S_T$ is convex and continuously differentiable, including
at the squared-hinge junction.  For any other feasible terminal action $w$,
its supporting inequality along the affine action direction gives
\[
 S_T(F(y,w,z))\ge S_T(F(y,v,z))
           +\partial_vS_T(F(y,v,z))(w-v).
\]
Use the same fixed feasible density $q$ at every $w$.  Weak duality
and its equality at $v$, followed by integration of this support, show
\[
 \rho_\alpha(S_T(F(y,w,\cdot)))
 \ge\rho_\alpha(S_T(F(y,v,\cdot)))
       +\mathbb E_P[q\,\partial_vS_T(F(y,v,Z))](w-v).
\]
Since $h''\ge2$, adding the original stage term proves the global
action minorant
\begin{equation}\label{eq:independentnodeminorant119}
 G_y(w)\ge G_y(v)+g(w-v)+(w-v)^2,
 \qquad 0\le w\le A(y).
\end{equation}
The action derivative inside the expectation is affine in $z$:
the quadratic parts of $S_T$ have affine gradients, while the
squared-hinge gradient is constant in $z$ along the zero-aggregate
shock.  Evaluations at $0,R$ therefore reconstruct this affine
derivative exactly.  No derivative of the optimal tail endpoints is used.

The minimum of $g\delta+\delta^2$ over the full feasible displacement
interval is attained at
$\delta=\operatorname{proj}_{[-v,A(y)-v]}(-g/2)$.
Minimizing the minorant gives the stated lower bound; evaluating the
feasible action $v$ gives the upper bound.  This reasoning applies if
$v=0$, $v=A(y)$, or an imperfect interior search iterate.  In
particular, a nonnegative support slope at zero certifies that endpoint
and a nonpositive support slope at capacity certifies that endpoint.
At an interior point, a negative support slope rules out a minimizer
to its left, and a positive support slope rules out one to its right.
Thus rational bisection is legitimate, but the support bound remains
sound after any finite number of steps.  The separately recorded check
$u_k-\ell_k\le2^{-20}$ is a width check; it is not the justification
of the lower bound.

\medskip\noindent
\emph{The continuous-risk outer enclosure.}
The original strong state convexity of $V_1$ implies convexity of
$f(z)=V_1(F(x,a,z))$, because the root transition is affine in $z$.
For every real threshold $\eta$, the function
$\phi_\eta(z)=(f(z)-\eta)_+$ is convex.  On the cell
$[z_{2j},z_{2j+2}]$, Jensen's inequality and convex secant domination
give, with cell length $\Delta=2R/M$,
\begin{align*}
 (\ell_{2j+1}-\eta)_+
 &\le\phi_\eta(z_{2j+1})
 \le\frac1\Delta\int_{z_{2j}}^{z_{2j+2}}\phi_\eta(z)\,dz,\\
 \frac1\Delta\int_{z_{2j}}^{z_{2j+2}}\phi_\eta(z)\,dz
 &\le\frac{\phi_\eta(z_{2j})+\phi_\eta(z_{2j+2})}{2}
 \le\frac{(u_{2j}-\eta)_++(u_{2j+2}-\eta)_+}{2}.
\end{align*}
Summing cells with weights $1/M$ and adding the threshold yields
\begin{align*}
 \eta+\frac1{Mt}\sum_{j=0}^{M-1}(\ell_{2j+1}-\eta)_+
 &\le\eta+t^{-1}\mathbb E_P(f-\eta)_+,\\
 \eta+t^{-1}\mathbb E_P(f-\eta)_+
 &\le\eta+\frac1{2Mt}\sum_{j=0}^{M-1}
             \{(u_{2j}-\eta)_++(u_{2j+2}-\eta)_+\}.
\end{align*}
These hold at every common threshold.  Taking infima over all real
thresholds preserves each inequality and proves
\[
 \rho_\alpha^{[M]}(\ell^{\mathrm{mid}})
 \le\rho_\alpha(f)
 \le\rho_\alpha^{[2M]}(u^{\mathrm{end}}).
\]
Neither the lower nor the upper node array has to be convex.  The
required convexity is that of the original optimal continuation along
the innovation.  A bound only on the ordinary integral of $f$ would
not prove this risk enclosure.  Adding $c$ and multiplying by the
positive $\beta$ gives~\eqref{eq:independentQ119}.

Finally, the threshold formula is monotone in its random argument,
and replacing its threshold by $\eta-c_0$ proves
$\rho_\alpha(w+c_0)=\rho_\alpha(w)+c_0$ for every constant
$c_0$.  Apply these facts to~\eqref{eq:independentterminalallowance119}
at the original successors to obtain
\[
 Q_x(a)\le J_x(a)\le Q_x(a)+\beta\varepsilon_1.
\]
Combining with the $Q$ bounds proves~\eqref{eq:independentJ119}.
The node-optimization uncertainty is already contained in $\ell_k,u_k$.
The distinct implemented-future allowance $\beta\varepsilon_1$ is
added exactly once at the root.  No convexity or differentiability of
the rounded deployed future controller is asserted.
\end{proof}

\subsection{Exact atom handling and finite policy comparisons}
For completeness, the equal-mass risk in the proposition can be computed
without floating cutoffs.  Sort $w_{(1)}\ge\cdots\ge w_{(n)}$,
put $r=n(1-\alpha)>0$, $k=\lfloor r\rfloor$, and
$\theta=r-k$.  Then
\begin{equation}\label{eq:independentatoms119}
 \rho_\alpha^{[n]}(w)
 =\frac{\displaystyle\sum_{i=1}^k w_{(i)}
                  +\theta w_{(k+1)}}{r},
\end{equation}
where the final term is omitted if $\theta=0$, including $k=n$.
Indeed, assign masses $\lambda_i\in[0,1]$ with
$\sum_i\lambda_i=r$.  The threshold inequality in the proof gives
$r^{-1}\sum_i\lambda_iw_i\le\rho_\alpha^{[n]}(w)$.
The assignment of unit mass to the $k$ largest values and mass
$\theta$ to the next value attains the threshold expression at the
cutoff value.  This proves the formula, including ties and split atoms.
At $\alpha=0$ it reduces to the ordinary mean; $\alpha=1$ is
outside this definition.  At the two frozen risks and the frozen
$M=64$, both atom counts are integral.  The implementation nevertheless
uses exact rational masses and preserves the repeated endpoints.

Suppose a fixed catalogue contains $n_0$ original states and each complete
policy $r$ has a valid interval $[L_{jr},U_{jr}]$ at state $j$.
Summing its pointwise inequalities gives the deterministic mean bound
\[
 \frac1{n_0}\sum_jL_{jr}
 \le\frac1{n_0}\sum_jJ_{jr}
 \le\frac1{n_0}\sum_jU_{jr}.
\]
For two policies $r,s$, use the interval
$[L_{jr}-U_{js},U_{jr}-L_{js}]$ for their difference at state $j$.
If they have the same exact action, the same original state and risk,
and the identical future controller, the two original defining objectives
are equal, and that state's difference is exactly zero.  Replacing its
subtraction interval by $[0,0]$ is then an identity, not an independence
assumption.  Summing these statewise difference intervals proves the
reported mean-contrast interval.  A strictly negative upper endpoint
certifies a lower cost for the left policy; a strictly positive lower
endpoint certifies a higher cost.  An interval containing zero remains
unresolved, except that $[0,0]$ certifies equality.  No statistical or
population inference follows from these deterministic sums.

In D, exact stored binary64 identities are converted to their precise
rational values.  The input binding checks the common state across every
action and service and the unchanged future-controller identity.  The
frozen catalogue has 144 states per risk, 983 distinct state/action/risk
keys, and 13,824 primary policy records.  The prescribed output includes
all 96 arm/service/risk means, all 56 unordered arm contrasts with the
fixed compiled-service labels, and all 240 within-arm service-pair
identity cells.  These are declared coverage requirements, not assertions
that execution succeeded.  No favorable sign is a pass condition.

\subsection{What interval compatibility does and does not establish}
Let $I$ be an independently justified interval for an original policy
value and $P_0$ a recorded service or fine-reference interval.  If they
are disjoint, both cannot enclose the same value: at least one enclosure
argument, arithmetic step, input identity, or model premise requires
investigation.  The case must be retained.  If they intersect, their
intersection alone does not establish correctness of $P_0$; for example,
$I=[0,2]$ and $P_0=[1,3]$ overlap although a true value $1/2$ lies
outside $P_0$.  The stronger inclusion $I\subseteq P_0$ would transfer
coverage from $I$ to $P_0$, conditional on validity of $I$.
The reverse inclusion does not.  D records overlap and both inclusion
directions separately and uses overlap only as a compatibility diagnostic.
It does not thereby verify every endpoint of the old 1,024-cell reference,
which shares the production terminal compiler.

The lower bound $L_Q$ in~\eqref{eq:independentQ119} concerns the
specified action.  It does not by itself provide a valid lower bound on
$\inf_aQ_x(a)$ and therefore creates no new global regret or
action-distance certificate.
The procedure leaves all original A/B/C policies, observations, timing
records, negative results, and qualification proofs unchanged.  Its full
actual work is additional author-side validation cost.  The subsequent
execution receipt, rather than this mathematical proposition, determines
whether every fixed case was actually completed and passed.
